% ID: SEL-D0025-OB-E01
% ESTADO: validada. S026, cotejada con examen y criterios.
\textbf{a)} $A$ es $2\times3$, $B$ es $3\times2$ y $C$ es $2\times2$.
\begin{itemize}
\item $AB^t$ es imposible porque sus dimensiones interiores son 3 y 2.
\item $B+3C$ es imposible porque los sumandos tienen distinto número de filas.
\item $CB^t$ es posible y da una matriz $2\times3$.
\item $AB+C$ es posible porque $AB$ y $C$ son $2\times2$.
\end{itemize}
\textbf{b)} $AB=\begin{pmatrix}1&-1\\-2&3\end{pmatrix}$ tiene
determinante 1. Despejamos multiplicando por su inversa a la izquierda:
\[
X=(AB)^{-1}C=
\begin{pmatrix}3&1\\2&1\end{pmatrix}\begin{pmatrix}1&1\\3&-2\end{pmatrix}
=\boxed{\begin{pmatrix}6&1\\5&0\end{pmatrix}}.
\]
Al multiplicar $AB$ por $X$ se recuperan las filas $(1,1)$ y $(3,-2)$.
